% ============================================================
% TAMBAHAN BAB III — Subbab baru di akhir rancangan-audit-record.tex
% Tempatkan setelah \subsubsection{Atribut Khusus Audit Record}
% dan sebelum penutup subsection tersebut.
% ============================================================

\subsubsection{Rancangan Skema Terperinci AuditEvent dan Enum Atribut}
\label{subsubsec:rancangan-skema-audit-event-detail}

Atribut umum yang telah diidentifikasi pada Subbab sebelumnya diturunkan
menjadi skema data terperinci yang terdiri atas dua \textit{struct} bersarang:
\texttt{AuditEvent} sebagai pembungkus tingkat atas dan \texttt{Event} sebagai
wadah seluruh atribut detail \textit{event}. Pemisahan ini dilakukan agar
atribut yang bersifat metadata pengiriman (\texttt{source\_component} dan
\texttt{source\_timestamp}) dapat dipisahkan secara struktural dari atribut
yang mendeskripsikan aktivitas (\texttt{actor\_id}, \texttt{action\_type},
\texttt{target\_object}, \texttt{outcome}, dan \texttt{details}).

\paragraph{Struct AuditEvent}

\texttt{AuditEvent} memuat tiga bidang tingkat atas: \texttt{source\_component}
bertipe \texttt{AuditSourceComponent} yang mengidentifikasi komponen pemangkit
\textit{event}, \texttt{source\_timestamp} bertipe \texttt{DateTime<Utc>} yang
mencatat waktu kejadian menurut jam komponen pengirim, dan \texttt{event} bertipe
\texttt{Event} yang memuat seluruh atribut aktivitas.

\paragraph{Dual Timestamp: source\_timestamp dan timestamp}

Skema ini secara sengaja membedakan dua nilai waktu. \texttt{source\_timestamp}
dicatat oleh \textit{event source} pada saat kejadian berlangsung dan disertakan
di dalam \textit{event} sebelum enkripsi, sehingga nilainya tidak dapat diubah
oleh ALS. \texttt{timestamp} pada \texttt{AuditRecord} dicatat oleh ALS pada
saat \textit{record} dibentuk dari antrian internal. Kedua nilai dapat berbeda
apabila terdapat latensi jaringan atau antrean yang signifikan. Penyertaan
\texttt{source\_timestamp} diperlukan agar urutan kejadian yang sesungguhnya
dapat direkonstruksi selama investigasi, terlepas dari urutan kedatangan
\textit{event} di sisi ALS.

\paragraph{Struct Event dan Atribut Berketik}

\texttt{Event} memuat tujuh bidang: \texttt{actor\_id} (identitas string aktor),
\texttt{actor\_type} bertipe \texttt{AuditActorType} (kategori aktor),
\texttt{target\_object} (identitas string objek target), \texttt{target\_object\_type}
bertipe \texttt{AuditTargetObjectType} (kategori objek target), \texttt{outcome}
bertipe \texttt{AuditOutcome}, \texttt{action\_type} bertipe \texttt{AuditActionType},
dan \texttt{details} bertipe \texttt{AuditEventDetails}.

Penggunaan \textit{enum} berketik pada \texttt{actor\_type},
\texttt{target\_object\_type}, \texttt{outcome}, dan \texttt{action\_type}
alih-alih representasi \texttt{String} bebas merupakan keputusan desain yang
disengaja. \textit{Enum} berketik memindahkan validasi nilai ke tahap
deserialisasi: nilai yang tidak dikenal langsung ditolak sebelum masuk ke
antrian pemrosesan, sehingga tidak diperlukan validasi terpisah di dalam logika
\textit{handler}. Pendekatan ini juga menjamin konsistensi format di seluruh
komponen karena hanya nilai yang telah didefinisikan pada \textit{enum} yang
dapat dikirimkan.

\paragraph{Enum AuditSourceComponent}

\texttt{AuditSourceComponent} mengkategorikan komponen pemangkit \textit{event}.
Keempat varian yang dirancang sesuai dengan \textit{event source} yang
diidentifikasi pada Subbab~\ref{sec:analisis-event-source} ditunjukkan pada
Tabel~\ref{tab:enum-source-component}.

\begin{table}[htbp]
\centering
\caption{Varian \texttt{AuditSourceComponent}}
\label{tab:enum-source-component}
\renewcommand{\arraystretch}{1.3}
\begin{tabular}{|p{4.5cm}|p{7.5cm}|}
\hline
\textbf{Varian} & \textbf{Komponen} \\
\hline
\texttt{HospitalClient} & Aplikasi \textit{Hospital Client} berbasis Tauri. \\
\hline
\texttt{PatientClient} & Aplikasi \textit{Patient Client} berbasis Tauri. \\
\hline
\texttt{MinistryClient} & Aplikasi \textit{Ministry Client} berbasis Tauri. \\
\hline
\texttt{ProxyReencryption} & PRE Server. \\
\hline
\end{tabular}
\end{table}

\paragraph{Enum AuditActionType}

\texttt{AuditActionType} mengkategorikan jenis tindakan yang dilakukan. Varian
dirancang berdasarkan operasi CRUD yang diperluas dengan tindakan khas sistem
dekMed. Nilai serialisasi JSON menggunakan format \texttt{SCREAMING\_SNAKE\_CASE}
agar mudah dibedakan dari nilai string bebas. Daftar varian yang dirancang
merupakan superset dari nilai-nilai yang telah didefinisikan pada
Tabel~\ref{tab:nilai-action-type}, dengan penambahan varian \texttt{SIGN\_IN}
untuk mencakup peristiwa autentikasi sesi yang terjadi pada aplikasi
\textit{client}.

\begin{table}[htbp]
\centering
\caption{Varian \texttt{AuditActionType}}
\label{tab:enum-action-type}
\renewcommand{\arraystretch}{1.3}
\begin{tabular}{|p{3cm}|p{9cm}|}
\hline
\textbf{Varian} & \textbf{Makna} \\
\hline
\texttt{CREATE} & Membuat atau menginisialisasi objek baru. \\
\hline
\texttt{READ} & Membaca data tanpa mengubah \textit{state}. \\
\hline
\texttt{UPDATE} & Memperbarui objek yang sudah ada. \\
\hline
\texttt{DELETE} & Menghapus atau mencabut objek. \\
\hline
\texttt{VALIDATE} & Memverifikasi keabsahan tanpa mengubah \textit{state}. \\
\hline
\texttt{EXECUTE} & Menjalankan proses yang telah disiapkan sebelumnya. \\
\hline
\texttt{SIGN\_IN} & Autentikasi sesi pengguna; hanya digunakan pada aplikasi \textit{client}. \\
\hline
\end{tabular}
\end{table}

\paragraph{Enum AuditActorType}

\texttt{AuditActorType} mengkategorikan jenis aktor yang memicu \textit{event}.
Varian mencakup seluruh peran yang terdefinisi pada sistem DecMed, ditambah
varian \texttt{Unknown} untuk skenario di mana identitas aktor belum dapat
ditentukan pada saat \textit{event} dibangkitkan (misalnya sebelum autentikasi
selesai). Daftar varian ditunjukkan pada Tabel~\ref{tab:enum-actor-type}.

\begin{table}[htbp]
\centering
\caption{Varian \texttt{AuditActorType}}
\label{tab:enum-actor-type}
\renewcommand{\arraystretch}{1.3}
\begin{tabular}{|p{4.5cm}|p{7.5cm}|}
\hline
\textbf{Varian} & \textbf{Jenis Aktor} \\
\hline
\texttt{AdministrativePersonnel} & Staf administrasi rumah sakit. \\
\hline
\texttt{MedicalPersonnel} & Tenaga medis (dokter, perawat, dan sejenisnya). \\
\hline
\texttt{Patient} & Pasien. \\
\hline
\texttt{Admin} & Administrator rumah sakit. \\
\hline
\texttt{Ministry} & Pengguna Kementerian Kesehatan. \\
\hline
\texttt{PREServer} & PRE Server yang bertindak sebagai aktor otomatis. \\
\hline
\texttt{Unknown} & Identitas aktor belum dapat ditentukan. \\
\hline
\end{tabular}
\end{table}

\paragraph{Enum AuditTargetObjectType}

\texttt{AuditTargetObjectType} mengkategorikan jenis objek yang menjadi target
aktivitas. Varian yang dirancang mencerminkan seluruh sumber daya yang dapat
diakses atau dimanipulasi dalam sistem DecMed, sebagaimana ditunjukkan pada
Tabel~\ref{tab:enum-target-object-type}.

\begin{table}[htbp]
\centering
\caption{Varian \texttt{AuditTargetObjectType}}
\label{tab:enum-target-object-type}
\renewcommand{\arraystretch}{1.3}
\begin{tabular}{|p{4.5cm}|p{7.5cm}|}
\hline
\textbf{Varian} & \textbf{Objek yang Dikategorikan} \\
\hline
\texttt{ActivationKey} & Kunci aktivasi akun pengguna. \\
\hline
\texttt{AccessCapability} & \textit{Capability} akses pada \textit{smart contract}. \\
\hline
\texttt{MedicalRecord} & Berkas rekam medis. \\
\hline
\texttt{MedicalRecordMetadata} & Metadata rekam medis yang tersimpan \textit{on-chain}. \\
\hline
\texttt{AccessDelegationQR} & QR \textit{code} delegasi akses. \\
\hline
\texttt{AdministrativeData} & Data administratif rumah sakit. \\
\hline
\texttt{Nonce} & \textit{Nonce} sesi yang tersimpan sementara. \\
\hline
\texttt{AccessKeys} & Pasangan kunci akses (termasuk kunci PRE). \\
\hline
\texttt{Transaction} & Transaksi IOTA yang dikirimkan ke jaringan. \\
\hline
\texttt{IPFSObject} & Objek yang tersimpan pada IPFS. \\
\hline
\texttt{GasReservation} & Reservasi \textit{gas} melalui Gas Station Server. \\
\hline
\texttt{PREEndPoint} & \textit{Endpoint} layanan PRE Server. \\
\hline
\texttt{OnChainData} & Data yang tersimpan langsung pada \textit{ledger} IOTA. \\
\hline
\end{tabular}
\end{table}